\section{A critical smoothing time for primes modulo four}\label{sec:four}

For the prime-only race put $B_h(t)=\PP_3(t)-\PP_1(t)$ at modulus $4$.
Let $\beta(s)=L(s,\chi_4)$ denote the Dirichlet beta function,
and let $\gamma_1$ be its least positive nontrivial zero ordinate.
The Gaussian width determines whether the class $1$ can lead at
arbitrarily large scales.

\begin{theorem}\label{th:transition}
There is a unique $h_*>0$ with the following transition. For every
$0<h<h_*$, each strict sign of $B_h$ occurs with positive limiting
density in $t=\log x$. For every $h>h_*$, $B_h(t)$ is bounded below
by a positive constant for all sufficiently large $t$. At $h=h_*$
the limiting density of $B_{h_*}(t)>0$ is one and
$\liminf_{t\to\infty}B_{h_*}(t)=0$. The transition satisfies
\[
                 \frac{\log 2}{\gamma_1^2}< h_* < \frac1{22}.
\]
\end{theorem}

Equivalently let \(R_h(X)\) denote the excess of class \(3\) over class \(1\) with unnormalized weights \((\log p)\exp[-(\log(p/X))^2/(4h)]\). The square completion gives
\[
 \frac{R_h(X)}{\sqrt{X}}=\sqrt{4\pi h}\,\e^{h/4} B_h(\log X+h).
\]
Hence, for \(h>h_*\), \(R_h(X)\ge c_h\sqrt X\) eventually for some \(c_h>0\). Below the threshold, each class has a lead of at least \(c_h\sqrt X\) on a set of positive logarithmic density, after choosing \(c_h>0\) sufficiently small.

Knapowski and Tur\'an proved eventual divergence of this weighted excess at the particular width \(h=1/4\) under GRH for \(\beta\) \cite{kt,bp,suzuki}. Neroli proves critical exponents and smoothing orders for exponential and iterated Mellin kernels \cite[Theorems 2.2 and 3.3]{neroli}. The positivity argument below shares the idea of averaging translates with that work. For logarithmic Gaussians the multiplier is exactly \(\exp(-h\gamma^2)\), which permits a certified enclosure for the transition and a bound on the density of reversals.

\begin{proof}
The alternating series gives $\beta(1/2)>0$. The squares of both
units modulo four are $1$, so removing the prime squares adds $1/2$
to the difference between classes $3$ and $1$. The prime-only
formula~\eqref{eq:asymp} gives
\[
 B_h(t)=f_h(\theta_t)+o(1),\qquad
 f_h(z)=\tfrac12+2\operatorname{Re}\sum_{\gamma>0}
                    m_\gamma e^{-h\gamma^2} z_\gamma .
\]
Here $m_\gamma$ is the multiplicity for $\beta$, and the compact
group $H$ is the closure of $\theta_t=(e^{i\gamma t})_{\gamma>0}$.
The sign of the sum follows from subtracting class $1$ from class
$3$, which takes the negative of the $\chi_4$ component of the
prime-weighted sum. Every orbit $z\theta_t$ is dense in $H$.
Every open neighborhood in $H$ has positive Haar measure and is
visited at arbitrarily large positive times. Consequently
$m(h):=\min_H f_h=\liminf_{t\to\infty} B_h(t)$.

The function $m(h)$ is continuous for $h>0$ by local uniform convergence.
For $u>0$ the heat semigroup gives the identity
\[
 f_{h+u}(z)=\int_{\mathbb R} k_u(v)f_h(z\theta_v)\,dv .
\]
The function $f_h$ is nonconstant, as its Fourier coefficient at
$\gamma_1$ is nonzero. Each orbit visits an open set where
$f_h>m(h)$, and $k_u$ is strictly positive. Thus
$f_{h+u}(z)>m(h)$ for each $z$, and compactness gives
$m(h+u)>m(h)$. As $h\to\infty$, $f_h\to1/2$ uniformly.

Suppose $f_h\ge0$. Its Haar integral is $1/2$, so
\[
 m_{\gamma_1}e^{-h\gamma_1^2}
 =\left|\int_H f_h(z)\overline{z_{\gamma_1}}\,dz\right|<\tfrac12.
\]
For strictness, equality in the integral triangle inequality would
force $z_{\gamma_1}$ to be constant almost everywhere on the set
where $f_h>0$, hence everywhere there by continuity and full support.
That set contains an orbit segment, on which the coordinate rotates
at nonzero speed. Since $m_{\gamma_1}\ge1$, this proves
$m((\log2)/\gamma_1^2)<0$.

It follows that $m(h_*)=0$ has exactly one solution. Below it, both
strict sign sets are open and nonempty on $H$, since the mean is
positive. The projection has an absolutely continuous limiting law,
so $f_h=0$ has measure zero and contains the boundaries of both sign sets. Equidistribution
and the $o(1)$ approximation give the stated densities, including
density one at $h_*$. Above $h_*$ the positive minimum gives a uniform
margin at all sufficiently large times. The explicit upper endpoint
is proved next.
\end{proof}

\paragraph{An upper bound from two values of the completed logarithmic derivative.}
Write $D(s)=\Xi'(s)/\Xi(s)$ for
$\Xi(s)=(4/\pi)^{(s+1)/2}\Gamma((s+1)/2)\beta(s)$.
This is an entire function of order one, even about $s=1/2$.
The paired Hadamard product for
$\Xi(1/2+u)/\Xi(1/2)$ has factors
$(1+u^2/\gamma^2)^{m_\gamma}$ for $\gamma>0$.
Any linear exponential prefactor is constant by evenness.
Logarithmic differentiation at $u=1,2$ yields the positive moment
\[
 S:=\sum_{\gamma}\frac{m_\gamma}{(1+\gamma^2)(4+\gamma^2)}
       =\frac{D(3/2)-D(5/2)/2}{3}<0.00168 ,
\]
where both signs of each ordinate are included. For $y\ge0$,
\[
 (1+y)(4+y)e^{-y/22}\le(y+5/2)^2 e^{-y/22}
 \le44^2 e^{-83/44}<294 .
\]
Hence $\sum_\gamma m_\gamma e^{-\gamma^2/22}<294(0.00168)<1/2$.
This gives $m(1/22)>0$ uniformly over all phases, including any
relations among the ordinates.

For reproducibility, write $D(s)=A(s)+\beta'(s)/\beta(s)$. At
$s=3/2,5/2$ respectively the two values of $2A(s)-\log(4/\pi)$
are $4-\gamma_{\rm E}-\pi/2-3\log2$ and
$4/3-\gamma_{\rm E}+\pi/2-3\log2$, by digamma recurrence and reflection.
The partial alternating sums through $n=2047$ bound $\beta(3/2)$
from below and $\beta'(3/2)$ from above. Those through $2049$ bound
$\beta(5/2)$ from above and $\beta'(5/2)$ from below. Indeed
$(\log n)n^{-s}$ decreases for $n\ge3$ at either value of $s$,
and the derivatives of the constant terms vanish.
These rationally enclosable sums give $D(3/2)<0.15427$ and
$D(5/2)>0.29847$. The included \texttt{bound.py} checks the
inequalities with integer square roots and the series
$\log((1+u)/(1-u))=2\sum_{j\ge0}u^{2j+1}/(2j+1)$ after reduction
to $0\le u\le1/3$.

\paragraph{Direct arithmetic check.}
At $h=0.04$ the table compares the prime-power difference
$Q=\QQ_3-\QQ_1$, the prime-only difference $B=B_h$, and
the truncated zero sum $Z=2\sum_{0<\gamma\le32}
e^{-h\gamma^2}\cos(\gamma t)$ over the ten roots detected by
\texttt{verify.py}. The two classes have identical subtracted
main terms. The program sums odd prime powers up to $24\,154\,953$.
The integral majorant for the omitted arithmetic tail is less than
$1.2\times10^{-21}$ throughout the table. The ordinates were found
by a numerical sign scan, which supplies an illustration of the
formula. The upper bound for $h_*$ uses the separate rational certificate.

\begin{center}
\begin{tabular}{r@{\qquad}r@{\qquad}r@{\qquad}r@{\qquad}r}
\toprule
$t$ & $Q$ & $Z$ & $B$ & $B-Z-1/2$\\
\midrule
 6 & $ 0.002964$ & $ 0.002829$ & $0.341832$ & $-0.160996$\\
 8 & $-0.209244$ & $-0.209251$ & $0.339891$ & $ 0.049142$\\
10 & $-0.418259$ & $-0.418259$ & $0.060718$ & $-0.021023$\\
12 & $-0.495944$ & $-0.495944$ & $0.002424$ & $-0.001632$\\
14 & $-0.389047$ & $-0.389047$ & $0.119904$ & $ 0.008951$\\
\bottomrule
\end{tabular}
\end{center}

The first computed ordinate is about $6.02095$, placing the
symbolic lower bound near $0.01912$. Although class $3$ leads at all
five sample times, Proposition~\ref{prop:enclosure} proves that both
classes lead with positive logarithmic density at this width.
